\documentclass[11pt]{article}
\usepackage[margin=1in]{geometry}
\usepackage{lmodern}
\usepackage[T1]{fontenc}
\usepackage{microtype}
\usepackage{amsmath,amssymb,mathtools}
\usepackage{booktabs}
\usepackage{enumitem}
\usepackage{xcolor}
\usepackage[most]{tcolorbox}
\usepackage[colorlinks=true,linkcolor=blue!50!black,urlcolor=blue!50!black]{hyperref}

\definecolor{keyblue}{RGB}{30,70,140}
\newtcolorbox{keyidea}[1][]{colback=keyblue!5,colframe=keyblue!70,
  fonttitle=\bfseries,title=#1,boxrule=0.6pt,arc=2pt,left=6pt,right=6pt}
\newtcolorbox{caution}[1][]{colback=orange!6,colframe=orange!70!black,
  fonttitle=\bfseries,title=#1,boxrule=0.6pt,arc=2pt,left=6pt,right=6pt}

\newcommand{\E}{\mathbb{E}}
\newcommand{\rxy}{\rho_{XY}}
\setlist{itemsep=2pt,topsep=4pt}

\title{\textbf{Optimal Life-Cycle Portfolio with Stochastic Labor Income}\\[4pt]
\large AMS 516 Project --- Problem 12: an explanatory overview}
\author{Deepak Singh}
\date{October 2026}

\begin{document}
\maketitle

\begin{abstract}
These notes explain Problem 12 of the AMS 516 project description: a worker who
chooses consumption and a stock/bond portfolio over a 30-year working life while
receiving a risky salary. We describe the economic story, walk through the
Hamilton--Jacobi--Bellman (HJB) equation term by term, derive the optimal controls,
build intuition through the concept of \emph{human capital}, show that the problem
reduces from two state variables to one, and list the analytical benchmarks and the
numerical difficulties relevant for a neural PDE solver.
\end{abstract}

\setcounter{tocdepth}{1}
\tableofcontents

%======================================================================
\section{The story in plain words}

Picture a 35-year-old with $T=30$ working years ahead. At each moment they
\begin{itemize}
  \item \textbf{earn a salary} $Y_t$, which grows on average but is hit by random,
        permanent shocks (raises, bad years, career luck);
  \item \textbf{hold financial wealth} $X_t$ (savings);
  \item \textbf{choose a consumption rate} $C_t \ge 0$;
  \item \textbf{choose a dollar amount $\pi_t$ in stocks}; the rest, $X_t-\pi_t$,
        sits in a safe bond paying $r$ (negative means borrowing).
\end{itemize}
They want a smooth, comfortable life without going broke. The central question is:

\begin{keyidea}[The question]
How should the fact that you \emph{own a future salary} change how you invest and
how much you spend today?
\end{keyidea}

This is the economics behind target-date retirement funds (``hold more stocks when
young, fewer when old''). Problem~12 lets us \emph{derive} that advice rather than
assume it.

%======================================================================
\section{The model}

\subsection{Wealth dynamics}
\begin{equation}\label{eq:wealth}
dX_t = \big(\underbrace{rX_t}_{\text{interest}}
      + \underbrace{\pi_t(\mu-r)}_{\text{risk premium}}
      - \underbrace{C_t}_{\text{spending}}
      + \underbrace{Y_t}_{\text{salary}}\big)\,dt
      + \underbrace{\sigma\pi_t\,dW^X_t}_{\text{stock risk}} .
\end{equation}
All wealth earns $r$; the amount in stocks earns an additional $\mu-r$ and carries
risk $\sigma\pi_t$. Because $\pi_t$ is in dollars, $\pi_t > X_t$ means leverage and
$\pi_t<0$ means shorting.

\subsection{Salary dynamics}
\begin{equation}\label{eq:income}
dY_t = \alpha Y_t\,dt + \beta Y_t\,dW^Y_t,
\qquad d\langle W^X, W^Y\rangle_t = \rxy\,dt .
\end{equation}
Salary is a geometric Brownian motion: $\alpha$ is average growth and $\beta$ the
volatility of \emph{permanent} shocks (a bad year lowers the whole future path).
$Y_t$ stays positive but can drift arbitrarily close to zero. The correlation
$\rxy$ measures how much wages move with the stock market.

\subsection{Objective}
\begin{equation}\label{eq:objective}
\max_{\{\pi_t,\,C_t\}}\;
\E\!\left[\int_0^T e^{-\rho t}\,U(C_t)\,dt + e^{-\rho T}\,U(X_T)\right],
\qquad U(c) = \frac{c^{1-\gamma}}{1-\gamma}.
\end{equation}
\begin{itemize}
  \item $\rho$ is impatience (not to be confused with the correlation $\rxy$;
        in code we will call them \texttt{rho} and \texttt{corr\_xy}).
  \item $\gamma$ is relative risk aversion; $1/\gamma$ is also the willingness to
        shift consumption across time.
  \item With $\gamma=3$, $U(c) = -\tfrac{1}{2c^2}$: always negative and
        $U(c)\to-\infty$ as $c\to0$. The agent fears very low consumption.
  \item $U(X_T)$ is a terminal (bequest / retirement-pot) value of wealth.
\end{itemize}

\subsection{Parameters}
The project description gives two parameter sets:

\begin{center}
\begin{tabular}{lcccccccccc}
\toprule
 & $r$ & $\mu$ & $\sigma$ & $\rho$ & $\alpha$ & $\beta$ & $\rxy$ & $\gamma$ & $X_0$ & $Y_0$ \\
\midrule
Set 1 & 0.02 & 0.06 & 0.18 & 0.03 & 0.01 & 0.10 & 0 (or 0.3) & 3 & 1 & 0.5 \\
Set 2 & 0.03 & 0.08 & 0.20 & 0.04 & 0.02 & 0.15 & 0.25 & 3 & 100 & 30 \\
\bottomrule
\end{tabular}
\end{center}
Both use $T=30$. Suggested domains: $x\in[0,500]$ (or $[-200,500]$ if borrowing is
allowed), $y\in[0,150]$.

%======================================================================
\section{The HJB equation, term by term}

Define the value function
\[
V(t,x,y) = \sup_{\pi,C}\,
\E\!\left[\int_t^T e^{-\rho(s-t)}U(C_s)\,ds + e^{-\rho(T-t)}U(X_T)
\;\Big|\; X_t=x,\;Y_t=y\right].
\]
It satisfies
\begin{equation}\label{eq:hjb}
\partial_t V - \rho V + \sup_{\pi\in\mathbb{R},\,C\ge0} \mathcal{H} = 0,
\qquad V(T,x,y) = U(x),
\end{equation}
with Hamiltonian
\begin{align*}
\mathcal{H} \;=\;& U(C) && \text{utility from consuming now}\\
&+ \big(rx + \pi(\mu-r) - C + y\big)\,V_x && \text{wealth drift $\times$ value of a dollar}\\
&+ \alpha y\,V_y && \text{salary drift $\times$ value of salary}\\
&+ \tfrac12\sigma^2\pi^2\,V_{xx} && \text{portfolio risk ($V_{xx}<0$, so it hurts)}\\
&+ \tfrac12\beta^2y^2\,V_{yy} && \text{salary risk}\\
&+ \rxy\,\sigma\beta\,\pi y\,V_{xy} && \text{covariance of portfolio and salary.}
\end{align*}

\paragraph{Reading the equation.}
By It\^o's lemma, $\mathcal{H}-U(C)$ is the expected rate of change of $V$ along the
state dynamics. The HJB therefore says: \emph{required return $\rho V$ = flow
utility + expected change in value}. The last term is where the interesting
economics lives: stock holdings interact with the risk already carried through
one's job. At $t=T$ working life ends and only wealth matters, hence
$V(T,x,y)=U(x)$, independent of $y$.

%======================================================================
\section{Optimal decisions}

\subsection{Consumption}
\[
\frac{\partial\mathcal{H}}{\partial C}=0
\;\Longrightarrow\;
U'(C^*) = V_x
\;\Longrightarrow\;
C^* = (V_x)^{-1/\gamma}.
\]
Consume until the utility of one more dollar spent equals the value of one more
dollar saved; $V_x$ is the shadow value of wealth.

\subsection{Portfolio}
\[
\frac{\partial\mathcal{H}}{\partial \pi}=0
\;\Longrightarrow\;
(\mu-r)V_x + \sigma^2\pi^* V_{xx} + \rxy\sigma\beta y\,V_{xy} = 0,
\]
which splits naturally into two pieces:
\begin{equation}\label{eq:pistar}
\pi^* \;=\;
\underbrace{\frac{\mu-r}{\sigma^2}\left(-\frac{V_x}{V_{xx}}\right)}_{\text{speculative (Merton) demand}}
\;+\;
\underbrace{\left(-\rxy\,\frac{\beta y}{\sigma}\,\frac{V_{xy}}{V_{xx}}\right)}_{\text{hedging demand}} .
\end{equation}
\begin{itemize}
  \item $-V_x/V_{xx}$ is \emph{risk tolerance}: exploit the risk premium in
        proportion to it.
  \item More salary typically lowers the marginal value of a dollar, so $V_{xy}<0$;
        with $V_{xx}<0$ and $\rxy>0$ the hedging demand is \textbf{negative}: your
        job already exposes you to the market, so you hold \emph{less} stock. A job
        that pays well in crashes ($\rxy<0$) would imply \emph{more} stock.
\end{itemize}

\subsection{The PDE we actually solve}
Substituting $C^*$ and $\pi^*$ back in, and using
$U(C^*) - C^*V_x = \frac{\gamma}{1-\gamma}V_x^{(\gamma-1)/\gamma}$ and
$\sup_\pi\{a\pi+\tfrac12 b\pi^2\} = -a^2/(2b)$ for $b<0$:
\begin{equation}\label{eq:hjb-reduced}
\boxed{
\partial_tV - \rho V + \frac{\gamma}{1-\gamma}V_x^{\frac{\gamma-1}{\gamma}}
+ (rx+y)V_x + \alpha yV_y + \tfrac12\beta^2y^2V_{yy}
- \frac{\big[(\mu-r)V_x + \rxy\sigma\beta y V_{xy}\big]^2}{2\sigma^2 V_{xx}} = 0 }
\end{equation}
This is a \textbf{fully nonlinear} PDE: $V_{xx}$ sits in a denominator and $V_x$ is
raised to a fractional power. It is well defined only while $V_x>0$ and $V_{xx}<0$,
which is the root of most numerical difficulties (Section~\ref{sec:difficulties}).

%======================================================================
\section{Key intuition: human capital}

\subsection{Riskless salary ($\beta=0$)}
If salary is certain, the future salary stream is worth exactly its present value,
like owning a bond:
\begin{equation}\label{eq:human-capital}
H(t,y) = \int_t^T y\,e^{(\alpha-r)(s-t)}\,ds
       = y\,\frac{1-e^{-(r-\alpha)(T-t)}}{r-\alpha}.
\end{equation}
The problem then has an exact solution: behave like Merton, but on
\textbf{total wealth} $x+H$:
\begin{equation}\label{eq:det-income}
V = \frac{f(t)^\gamma\,(x+H)^{1-\gamma}}{1-\gamma},\qquad
\pi^* = \frac{\mu-r}{\gamma\sigma^2}\,(x+H),\qquad
C^* = \frac{x+H}{f(t)},
\end{equation}
where
\[
f(t) = \frac{1-e^{-\nu(T-t)}}{\nu} + e^{-\nu(T-t)},
\qquad
\nu = \frac{1}{\gamma}\left[\rho - (1-\gamma)\left(r + \frac{(\mu-r)^2}{2\gamma\sigma^2}\right)\right].
\]
(With $H\equiv0$ this is the classical Merton problem with consumption and bequest.)

\subsection{What it implies (parameter set 2, $t=0$)}
\begin{center}
\begin{tabular}{lr}
\toprule
Quantity & Value \\
\midrule
Human capital $H = 30\times 25.9$ & $\approx 778$ \quad(vs.\ savings of 100) \\
Merton fraction $(\mu-r)/(\gamma\sigma^2) = 0.05/0.12$ & $\approx 0.417$ \\
$\nu$,\; $f(0)$ & $\approx 0.0403$,\; $\approx 17.7$ \\
Stock holding $\pi^* = 0.417\times 878$ & $\approx 366$ \quad($3.7\times$ financial wealth) \\
Consumption $C^* = 878/17.7$ & $\approx 49.5$ \quad(salary is only 30) \\
\bottomrule
\end{tabular}
\end{center}

\begin{keyidea}[The life-cycle glide path]
A young worker should invest aggressively relative to savings (the large implicit
``bond'' $H$ provides safety) and consume more than they earn (borrowing against
future salary). As they age, $H$ shrinks to $0$ at $T$, and the stock share of
savings $\pi^*/x = \frac{\mu-r}{\gamma\sigma^2}(1 + H/x)$ falls toward the plain
Merton fraction. \textbf{That decline is the life-cycle glide path.}
\end{keyidea}

\subsection{Risky salary ($\beta>0$): the interesting case}
Now human capital is risky: partly market-like (through $\rxy$) and partly purely
personal risk that cannot be traded away. The market is \textbf{incomplete} and
there is no closed form. Three effects appear:
\begin{enumerate}
  \item \textbf{Hedging.} The market-correlated part makes human capital look a bit
        like stock, so the agent holds less stock (negative hedging demand
        in~\eqref{eq:pistar}).
  \item \textbf{Precautionary saving.} The uninsurable part makes the agent consume
        less and save more as a buffer.
  \item \textbf{No borrowing against salary.} Salary can collapse toward $0$ and
        $U\to-\infty$ at zero consumption, so the agent effectively will not borrow
        against it: $x\ge0$ becomes a self-imposed constraint. Near $x=0$, behaviour
        changes sharply ($\pi^*\to0$, consumption held near or below salary).
\end{enumerate}
This case is exactly what the neural PDE solver is for, and where the economics is
richest.

%======================================================================
\section{Dimension reduction: only the wealth-to-income ratio matters}

The dynamics \eqref{eq:wealth}--\eqref{eq:income} are linear in $(x,y,\pi,C)$ and
CRRA utility is homogeneous: $U(\lambda c)=\lambda^{1-\gamma}U(c)$. Scaling wealth,
salary and both controls by $\lambda>0$ therefore scales the whole problem, so
\[
V(t,\lambda x,\lambda y) = \lambda^{1-\gamma}V(t,x,y)
\quad\Longrightarrow\quad
\boxed{V(t,x,y) = y^{1-\gamma}\,v(t,z),\qquad z = \frac{x}{y}.}
\]
Consequences:
\begin{itemize}
  \item The 2-D (plus time) HJB collapses to a 1-D (plus time) PDE for $v(t,z)$, with
        controls per unit of salary, $c=C/y$ and $p=\pi/y$. This gives cheap,
        accurate finite-difference \textbf{ground truth} for the incomplete-market case.
  \item Only $z_0=X_0/Y_0$ matters, not the absolute scale: Set~1 has $z_0=2$ and
        Set~2 has $z_0\approx3.3$; the factor of 100 between them is irrelevant.
  \item $z$ (the wealth-to-income ratio) is the variable economists use to present
        policies in this literature.
\end{itemize}
Deriving the reduced PDE for $v(t,z)$ via the chain rule is a natural next step.

%======================================================================
\section{Benchmarks for validation}
\begin{enumerate}
  \item \textbf{No income} ($y=0$): classical Merton solution, \eqref{eq:det-income}
        with $H=0$. Also the exact boundary value at $y=0$, since salary that hits
        zero stays there.
  \item \textbf{Deterministic income} ($\beta=0$): human-capital solution
        \eqref{eq:human-capital}--\eqref{eq:det-income}.
  \item \textbf{Perfectly correlated income} ($\rxy=\pm1$): the market is complete;
        human capital is priced under the risk-neutral measure and the solution is
        again Merton on total wealth plus an explicit hedging term.
  \item \textbf{General case}: finite-difference / policy-iteration solution of the
        reduced 1-D PDE in $z$.
  \item \textbf{Monte Carlo policy evaluation}: simulate
        \eqref{eq:wealth}--\eqref{eq:income} under any learned policy and estimate
        its expected utility. No policy can beat the optimum, so this is a
        method-independent check.
\end{enumerate}

\begin{caution}[Note on the project description]
The remark lists ``$\sigma=0$'' as a special case with an analytical benchmark.
This almost certainly means $\beta=0$ (deterministic income): a riskless stock with
$\mu\ne r$ would be an arbitrage.
\end{caution}

%======================================================================
\section{Difficulties and modelling issues}\label{sec:difficulties}
\begin{enumerate}
  \item \textbf{Huge range of values.} $V<0$ and $V\to-\infty$ as $x\to0$. Fitting
        $V$ directly lets the region near $x=0$ dominate the loss. Better: learn a
        rescaled quantity such as the certainty-equivalent wealth
        $\big((1-\gamma)V\big)^{1/(1-\gamma)}$, or build the Merton form into the
        network.
  \item \textbf{Fragile controls.} $C^*$ needs $V_x>0$ and $\pi^*$ needs $V_{xx}<0$.
        If a network violates either during training, \eqref{eq:hjb-reduced} produces
        NaNs or nonsense. Remedies: architectures that enforce monotonicity and
        concavity, separate policy networks (actor--critic / policy iteration), or
        Deep BSDE methods.
  \item \textbf{Leverage near small $x$.} Optimal $\pi^*/x$ is large when $H\gg x$,
        giving steep features near $x=0$; collocation points should concentrate there.
  \item \textbf{The $x=0$ boundary} is a genuine economic constraint, not merely the
        edge of the grid. Our borrowing assumption must be stated explicitly.
  \item \textbf{The $y=0$ boundary} is a gift: the exact Merton solution.
  \item \textbf{Terminal condition quirk.} Since $V(T,x,y)=U(x)$, the consumption rule
        gives $C^*(T)=x$ exactly: the agent reaches $T$ with only about one year's
        consumption saved. For a retirement-savings story, a natural extension is a
        bequest weight $b\,U(X_T)$ or an explicit retirement phase.
  \item \textbf{Long horizon.} $T=30$ years means value and policies change
        substantially over time; time sampling and normalisation ($t/T$) matter.
\end{enumerate}

%======================================================================
\section{Project in one sentence}

\begin{keyidea}
Solve a fully nonlinear two-dimensional HJB with neural networks, validate it
against the Merton and human-capital formulas and a one-dimensional
finite-difference solver, then use it to show how risky, partly uninsurable labor
income shapes the life-cycle glide path, precautionary saving, and the
stock--salary hedge.
\end{keyidea}

\subsection*{Further reading}
\begin{itemize}
  \item Merton, R.\,C.\ (1969, 1971). Lifetime portfolio selection under uncertainty;
        Optimum consumption and portfolio rules in a continuous-time model.
  \item Bodie, Z., Merton, R.\,C., \& Samuelson, W.\ (1992). Labor supply flexibility
        and portfolio choice in a life cycle model.
  \item Viceira, L.\ (2001). Optimal portfolio choice for long-horizon investors with
        nontradable labor income.
  \item Cocco, J., Gomes, F., \& Maenhout, P.\ (2005). Consumption and portfolio choice
        over the life cycle.
  \item Achdou, Y., Han, J., Lasry, J.-M., Lions, P.-L., \& Moll, B.\ (2022). Income and
        wealth distribution in macroeconomics: a continuous-time approach.
  \item Sirignano, J., \& Spiliopoulos, K.\ (2018). DGM: a deep learning algorithm for
        solving partial differential equations.
\end{itemize}

\end{document}
